% 从已验证能力到合法派生 — 面向 Agent 预研的形式方法教材草案
% 编译：bash ../../scripts/xelatex-clean.sh 从已验证能力到合法派生-面向Agent预研的形式方法教材草案.tex
% 维护：单文件持续改写（无日期后缀）；工作分支 research/formal-lang-dag
\documentclass[UTF8,a4paper,11pt]{ctexart}
\usepackage{amsmath,amssymb}
\usepackage{booktabs}
\usepackage{geometry}
\usepackage{hyperref}
\usepackage{longtable}
\usepackage{array}
\usepackage{listings}
\usepackage{enumitem}
\usepackage[most]{tcolorbox}
\usepackage{tikz}
\usetikzlibrary{arrows.meta,positioning,fit,shapes.geometric}
\geometry{margin=2.2cm}
\newcolumntype{L}[1]{>{\raggedright\arraybackslash}p{#1}}
\hypersetup{colorlinks=true,linkcolor=blue,urlcolor=blue}
\lstset{basicstyle=\ttfamily\small,breaklines=true,frame=none,columns=fullflexible}
\setlist[itemize]{leftmargin=1.6em}
\setlist[enumerate]{leftmargin=1.8em}
% FIPS 算法框：黑字可读 + 参数/Input/Output + 与标准一致的步骤编号
\newenvironment{fipsalgo}[1]{%
  \begin{tcolorbox}[
    enhanced, breakable,
    colback=gray!4, colframe=gray!55, coltext=black,
    boxrule=0.5pt, arc=2pt,
    left=6pt, right=6pt, top=4pt, bottom=4pt,
    title={#1},
    fonttitle=\bfseries\color{black},
    coltitle=black,
    attach boxed title to top left={yshift=-2mm,xshift=4mm},
    boxed title style={colback=white, colframe=gray!55, boxrule=0.4pt, coltitle=black}
  ]
  \small\color{black}
}{%
  \end{tcolorbox}
}
% 算法框内：参数集 / Input / Output 行（FIPS 原文措辞）
\newcommand{\fipsioparam}[1]{\par\noindent\textbf{参数集}（\texttt{ml\_kem\_1024}）：#1\par\vspace{2pt}}
\newcommand{\fipsioin}[1]{\noindent\textbf{Input：}#1\par}
\newcommand{\fipsioout}[1]{\noindent\textbf{Output：}#1\par\vspace{4pt}}
\newenvironment{fipssteps}{%
  \begin{enumerate}[leftmargin=2.2em, label=\textbf{\arabic*:}, itemsep=2pt, topsep=2pt]
}{%
  \end{enumerate}
}

\title{从已验证能力到合法派生\\
{\large 面向 Agent 的预研开发形式方法\\（教材式草案·持续修订）}}
\author{ascendc 工程 · 专题分支 \texttt{research/formal-lang-dag}}
\date{}

\begin{document}
\maketitle

\begin{abstract}
\textbf{愿景}：把「层层递进、先验证底层再拼高层」的工程直觉，提升为一套可教学、可门禁、可对照实验的形式方法——
使 Agent（以及人类协作者）在复杂预研任务中，\textbf{只沿已验证能力及其合法组合}生长实现，
从机制上压缩跳层发明与幻觉空间；最终服务一般预研开发，而非仅服务某一密码标准的算子清单。

\textbf{写法}：仿教材而非论文——先从特例讲故事，再抽象数学对象，再复盘与前瞻；
证明与完备性后置，允许边做边补。\textbf{地位}：\texttt{docs/research/} 调研草稿，非定稿、非交付验收依据。

\textbf{压力测试床}：ML-KEM-1024 链上的 \texttt{pke.keygen} / \texttt{kem.keygen} / \texttt{kem.encaps}，
并以 \texttt{kem.decaps} 为理论前瞻试金石；\texttt{*-correctness-*} 作为「只求正确、不计优化」的对照组标本。
\end{abstract}

\tableofcontents
\newpage

% ============================================================
\section{愿景、读者与本文怎么读}

\subsection{愿景（一句话）}
我们想表达的是：预研开发\textbf{不应当}变成「对着算法说明书自由发挥写核」。
我们应当只在\textbf{已经拿到证书}的能力之上，按规则组合出更复杂的任务；
静态 DAG 只是「引理库谁依赖谁」的一张投影图，真正的作业过程还要另说。

\subsection{我们希望最终得到什么}
\begin{enumerate}
  \item 一套\textbf{教材式}叙述：新人 / 新 Agent 能从特例读懂「图从哪来、树怎么长」；
  \item 一套\textbf{可操作门禁}：任务启动必须先交依赖闭包与缺项，再写码；
  \item 一套\textbf{对照实验协议}：方法论指导下的实现 vs correctness 标本 vs 人工指导路径；
  \item （成型后）可解释的接受/拒绝轨迹，以及必要的证明——\textbf{不}作为本稿前置条件。
\end{enumerate}

\subsection{研究路线：从特殊到一般}
\begin{center}
\begin{tabular}{@{}L{2.8cm}L{6.2cm}L{5cm}@{}}
  \toprule
  阶段 & 做什么 & 刻意后置 \\
  \midrule
  特例引出 & 工程故事 $\to$ 图与语法树 & 完备公式 \\
  套公式 & 用数学\textbf{复述}已见现象 & 证明 \\
  改公式 & 缺口反过来改理论 & 为漂亮删工程事实 \\
  套别例 & Encaps 复盘；换域；Decaps 前瞻 & 一次塞太多缺口 \\
  迭代方法论 & 门禁与流程收敛 & 急进 Rule/Skill \\
  成型后 & 证明、可解释性 & 过早锁死抽象 \\
  \bottomrule
\end{tabular}
\end{center}

\subsection{阅读地图（教材顺序）}
\begin{enumerate}
  \item 第~\ref{sec:hook}~章：\textbf{完整例子}（KEM.KeyGen：算法原文、我们有什么、怎样做出来、引出课题）；
  \item 第~\ref{sec:mathbg}~章：\textbf{方法论数学背景}（图、闭包、形式语言、自动机直觉——\textbf{不含}密码/NTT 代数）；
  \item 第~\ref{sec:math}--\ref{sec:gates}~章：\textbf{装进工程对象与门禁}（节点/边/证书；Agent 工作流）；
  \item 第~\ref{sec:replay}--\ref{sec:forward}~章：\textbf{复盘与前瞻}（用理论看实现）；
  \item 第~\ref{sec:correctness}~章：\textbf{correctness 对照组}；
  \item 第~\ref{sec:open}~章：未决议事项与维护约定。
\end{enumerate}

讨论过程的流水账见同目录 \texttt{形式语言DAG预研讨论纪要.md}（单文件持续刷新）。

% ============================================================
\section{一个完整例子：我们怎样做出 ML-KEM-1024 的 KEM.KeyGen}
\label{sec:hook}

本章用一个\textbf{已经做完}的真实任务当例子。
我们先说清楚「要算什么」，再贴 FIPS~203 里的算法原文，
然后交代我们手里已经有什么、前期做过哪些实验、最后怎样拼出 KEM.KeyGen。
例子讲完之后，我们再\textbf{提问}——这些问题应当能把读者自然带到后文的研究课题里。

\subsection{我们要交付什么}

我们的目标是：在昇腾 AI Core 上用 AscendC 实现
\textbf{FIPS~203 ML-KEM-1024 的密钥生成}（\texttt{ML-KEM.KeyGen}，即 Algorithm~19）。

读者可以把交付物理解成下面这张黑盒表（细节以 stable 目录 customspec 为准）：

\begin{center}
\begin{tabular}{@{}lll@{}}
  \toprule
  方向 & 内容 & 说明 \\
  \midrule
  输入 & \texttt{seed\_d.bin}（4\,B） & Host 写入；设备侧派生 $d$、$z$ \\
  输入 & NTT LUT（静态） & 与 PKE KeyGen 同表 \\
  输出 & \texttt{ek\_kem.bin}（1568\,B） & 与 \texttt{ek\_PKE} 相同 \\
  输出 & \texttt{dk\_kem.bin}（3168\,B） & 与 liboqs 展开布局一致 \\
  \bottomrule
\end{tabular}
\end{center}

我们\textbf{不}把 $d$、$z$、中间 $\hat{A}$、$\hat{s}$ 等当作交付文件落盘。
验收方式是：\texttt{run.sh} 跑 CPU 孪生 + SIM，输出与 golden / liboqs 对拍一致。

\subsection{FIPS~203 原文：三层算法怎样套在一起}

标准把 KEM 密钥生成拆成三层。
下面给出\textbf{与 FIPS~203 原文一致的参数、Input/Output 与步骤编号}。
正文引用 NIST FIPS~203 Algorithm~13 / 16 / 19；本文实例固定为 \texttt{ml\_kem\_1024}（$k{=}4$ 时 $\lvert\mathit{ek}\rvert{=}1568$\,B，$\lvert\mathit{dk}\rvert{=}3168$\,B）。

\begin{fipsalgo}{Algorithm 19 \texttt{ML-KEM.KeyGen()}（FIPS~203 §7.1）}
\fipsioparam{$n{=}256$，$q{=}3329$，$k{=}4$，$\eta_1{=}\eta_2{=}2$，$d_u{=}11$，$d_v{=}5$。}
\fipsioin{（无显式输入；$d$、$z$ 在算法内由 RBG 采样，见步骤 1--2。）}
\fipsioout{封装密钥 $\mathit{ek} \in \mathbb{B}^{384k+32}$；\\
           解封密钥 $\mathit{dk} \in \mathbb{B}^{768k+96}$。}
\begin{fipssteps}
  \item $d \leftarrow\$ \mathbb{B}^{32}$
  \item $z \leftarrow\$ \mathbb{B}^{32}$
  \item \textbf{if} $d = \mathrm{NULL}$ \textbf{or} $z = \mathrm{NULL}$ \textbf{then}
  \item \quad\textbf{return} $\bot$ \hfill{\footnotesize RBG 失败指示}
  \item \textbf{end if}
  \item $(\mathit{ek}, \mathit{dk}) \leftarrow \texttt{ML-KEM.KeyGen\_internal}(d, z)$
  \item \textbf{return} $(\mathit{ek}, \mathit{dk})$
\end{fipssteps}
\end{fipsalgo}

\begin{fipsalgo}{Algorithm 16 \texttt{ML-KEM.KeyGen\_internal}$(d, z)$（FIPS~203 §6.1）}
\fipsioparam{$n{=}256$，$q{=}3329$，$k{=}4$，$\eta_1{=}\eta_2{=}2$，$d_u{=}11$，$d_v{=}5$。}
\fipsioin{随机种子 $d \in \mathbb{B}^{32}$；随机种子 $z \in \mathbb{B}^{32}$。}
\fipsioout{封装密钥 $\mathit{ek} \in \mathbb{B}^{384k+32}$；\\
           解封密钥 $\mathit{dk} \in \mathbb{B}^{768k+96}$。}
\begin{fipssteps}
  \item $(\mathit{ek}_{\mathrm{PKE}}, \mathit{dk}_{\mathrm{PKE}}) \leftarrow \texttt{K-PKE.KeyGen}(d)$
  \item $\mathit{ek} \leftarrow \mathit{ek}_{\mathrm{PKE}}$
  \item $\mathit{dk} \leftarrow \mathit{dk}_{\mathrm{PKE}} \Vert \mathit{ek} \Vert H(\mathit{ek}) \Vert z$
  \item \textbf{return} $(\mathit{ek}, \mathit{dk})$
\end{fipssteps}
\end{fipsalgo}

\noindent 其中 $H = \mathrm{SHA3\text{-}256}$（FIPS~203 记号 $H$）。

\begin{fipsalgo}{Algorithm 13 \texttt{K-PKE.KeyGen}$(d)$（FIPS~203 §5.1）}
\fipsioparam{$n{=}256$，$q{=}3329$，$k{=}4$，$\eta_1{=}\eta_2{=}2$，$d_u{=}11$，$d_v{=}5$。}
\fipsioin{随机种子 $d \in \mathbb{B}^{32}$。}
\fipsioout{加密密钥 $\mathit{ek}_{\mathrm{PKE}} \in \mathbb{B}^{384k+32}$；\\
           解密密钥 $\mathit{dk}_{\mathrm{PKE}} \in \mathbb{B}^{384k}$。}
\begin{fipssteps}
  \item $(\rho, \sigma) \leftarrow G(d \Vert k)$ \hfill{\footnotesize 域分离：字节 33 为模块维 $k$}
  \item $N \leftarrow 0$
  \item \textbf{for} $i \leftarrow 0$ \textbf{to} $k-1$ \textbf{do}
  \item \quad\textbf{for} $j \leftarrow 0$ \textbf{to} $k-1$ \textbf{do}
  \item \quad\quad$\hat{A}[i,j] \leftarrow \texttt{SampleNTT}(\rho \Vert j \Vert i)$
  \item \quad\textbf{end for}
  \item \textbf{end for}
  \item \textbf{for} $i \leftarrow 0$ \textbf{to} $k-1$ \textbf{do}
  \item \quad$s[i] \leftarrow \texttt{SamplePolyCBD}_{\eta_1}(\texttt{PRF}_{\eta_1}(\sigma, N))$;\quad $N \leftarrow N+1$
  \item \textbf{end for}
  \item \textbf{for} $i \leftarrow 0$ \textbf{to} $k-1$ \textbf{do}
  \item \quad$e[i] \leftarrow \texttt{SamplePolyCBD}_{\eta_1}(\texttt{PRF}_{\eta_1}(\sigma, N))$;\quad $N \leftarrow N+1$
  \item \textbf{end for}
  \item $\hat{s} \leftarrow \texttt{NTT}(s)$
  \item $\hat{e} \leftarrow \texttt{NTT}(e)$
  \item $\hat{t} \leftarrow \hat{A} \circ \hat{s} + \hat{e}$ \hfill{\footnotesize NTT 域矩阵–向量乘加}
  \item $\mathit{ek}_{\mathrm{PKE}} \leftarrow \texttt{ByteEncode}_{12}(\hat{t}) \Vert \rho$
  \item $\mathit{dk}_{\mathrm{PKE}} \leftarrow \texttt{ByteEncode}_{12}(\hat{s})$
  \item \textbf{return} $(\mathit{ek}_{\mathrm{PKE}}, \mathit{dk}_{\mathrm{PKE}})$
\end{fipssteps}
\end{fipsalgo}

\textbf{读者只需抓住一条结构链}：
$\text{Alg.19} \Rightarrow \text{Alg.16} \Rightarrow \text{Alg.13}$。
Alg.19 比 Alg.13 \textbf{多出来的}主要是步骤 1--2 的 $z$、Alg.16 步骤 3 的拼接——
\textbf{并不是}把 NTT 或矩阵乘从头再写一遍。

\paragraph{本仓工程差异（验收已锁定，写在算法旁便于对照）}
\begin{itemize}
  \item 我们可复现跑分时，Host 只提供 4\,B \texttt{seed\_d}；$d$ 与 $z$ 在设备 UB 内派生（对应 Alg.19 步骤 1--2 的确定性扩展）。
  \item 我们对拍 liboqs 时，$\mathit{dk}$ 为 3168\,B 展开式
        $\mathit{dk\_pke}\Vert\mathit{ek}\Vert H\Vert z$，与 FIPS 代数最小形态一致，但字节布局与偏移以 customspec 为准。
\end{itemize}

\subsection{动手之前，我们手里已经有什么}

到做 KEM.KeyGen 之前，本仓库并不是从一张白纸开始。
我们已经在 AscendC 上做过大量单点与分段实验，并且多数有 CPU+SIM 对拍记录。
下面这张\textbf{资产清单}把「当时手里有什么」写清楚；文字说明与表格对照阅读。

\paragraph{四类家底（文字）}
\begin{enumerate}
  \item \textbf{平台层（P0）}：我们验证过 DataCopy、TPipe/TQue、MIX 核同步、MMAD tiling 等。
  \item \textbf{密码原语层（P1）}：我们验证过 NTT（poly-batch）、basemul、SHAKE/SHA3、CBD、ByteEncode$_{12}$、SampleNTT 等。
  \item \textbf{PKE 全链（P2）}：我们已有交付级 stable PKE KeyGen（2 launch，KAT 通过）。
  \item \textbf{对照标本}：我们保留 \texttt{*-correctness-*}，只求可执行与字节正确，不计优化。
\end{enumerate}

\paragraph{资产清单（表格；锚点以活跃 INDEX 为准）}
{\small
\begin{longtable}{@{}L{1.1cm}L{3.6cm}L{4.8cm}L{2.2cm}L{2.0cm}@{}}
  \toprule
  层 & 能力（简述） & 仓库锚点（示例） & 证书 & 备注 \\
  \midrule
  P0 & DataCopy / 搬运 & \texttt{docs/notes/ascendc-DataCopy…} & 文档+探针 & 平台知识库 \\
  P0 & MMAD + tiling & \texttt{pass-int8-matmul-cube-…} 等 & CPU+SIM & 闭合条件见 notes \\
  P0 & TPipe / 细同步 & KeyGen prep 技术总结 & CPU+SIM & 翻车高发区 \\
  \midrule
  P1 & NTT 三段式 poly-batch & \texttt{pass-fix-…-vec-k4-v2} & CPU+SIM & ML-KEM 活跃基线 \\
  P1 & SampleNTT / Alg.7 & \texttt{pass-fix-f203-alg7-sample-ntt-k4} & CPU+SIM & \\
  P1 & CBD $\eta{=}2$ & \texttt{pass-fix-f203-alg8-cbd-eta2-k4} & CPU+SIM & \\
  P1 & ByteEncode$_{12}$ & \texttt{pass-fix-…-byteencode12-vec-k4} & CPU+SIM & \\
  \midrule
  P2 & Alg.13 行 3--7 $\hat{A}$ & \texttt{pass-fix-f203-alg13-lines3-7-a-hat-k4} & CPU+SIM & \\
  P2 & Alg.13 行 8--15 $s,e$ & \texttt{pass-fix-f203-alg13-lines8-15-se-k4} & CPU+SIM & \\
  P2 & Alg.13 device 全链 & \texttt{pass-fix-f203-alg13-device-keygen-k4} & CPU+SIM & $\approx$542k tick \\
  P2 & \textbf{stable PKE KeyGen} & \texttt{stable-fips203-mlkem-pke-keygen-k4} & 交付+KAT & \textbf{KEM 直接依赖} \\
  \midrule
  P3 & KEM KeyGen correctness & \texttt{fix-f203-alg19-kem-keygen-correctness-k4} & CPU+SIM+liboqs & 对照标本 \\
  P3 & KEM KeyGen device & \texttt{pass-fix-f203-alg19-kem-keygen-device-k4} & CPU+SIM & $\approx$713k tick \\
  P3 & \textbf{stable KEM KeyGen} & \texttt{stable-fips203-mlkem-kem-keygen-k4} & 交付+KAT & 本例终点 \\
  \bottomrule
\end{longtable}
}

当我们决定做 KEM.KeyGen 时，
\textbf{我们并不是让 Agent 从 FIPS 伪代码一行行翻译成 AscendC}，
而是心里已经有上表这张「哪些积木是绿的」的清单。

\subsection{我们怎样探出来：工程路径与数据依赖}

这张清单不是一天长出来的。
本节用两张图把同一件事说清楚：
\textbf{第一张}是仓库时间线（先长 PKE，再接 KEM 尾缝）；
\textbf{第二张}是 FIPS Alg.19$\to$16$\to$13 的数据依赖（与 §2.2 步骤编号对照）。

\paragraph{图一：工程探路（PKE 线 $\to$ KEM 线）}
\begin{center}
\begin{tikzpicture}[
  node distance=6mm and 8mm,
  box/.style={draw, rounded corners=2pt, align=center, font=\small,
              minimum width=2.6cm, minimum height=0.85cm},
  arr/.style={-{Stealth[length=2mm]}, thick}
]
  % PKE 路径
  \node[box, fill=blue!8] (v2) {P1：NTT+UB\\vec-k4-v2};
  \node[box, fill=blue!8, right=of v2] (ahat) {Alg.13 行 3--7\\$\hat{A}$ 探针};
  \node[box, fill=blue!8, right=of ahat] (se) {Alg.13 行 8--15\\$s,e$ 探针};
  \node[box, fill=blue!12, right=of se] (pke) {P2：stable\\PKE KeyGen};
  \draw[arr] (v2) -- (ahat) -- (se) -- (pke);

  % KEM 路径
  \node[box, fill=orange!10, below=12mm of pke] (corr) {correctness\\（3 launch 标本）};
  \node[box, fill=orange!14, left=of corr] (dev) {device\\2 launch+尾缝};
  \node[box, fill=orange!18, left=of dev] (sk) {stable\\KEM KeyGen};
  \draw[arr] (pke) -- (corr);
  \draw[arr] (pke) -- (dev);
  \draw[arr] (corr) -- node[above, font=\footnotesize]{oracle} (dev);
  \draw[arr] (dev) -- (sk);

  \node[font=\footnotesize, align=left, anchor=north west] at ([yshift=-2mm]v2.south west) {
    \textbf{PKE 线}：先原语 $\to$ 分段 $\to$ 集成晋级\\
    \textbf{KEM 线}：复用 PKE $\to$ Alg.16 尾缝 $\to$ 去 vendor 晋级
  };
\end{tikzpicture}
\end{center}

\begin{itemize}
  \item \textbf{PKE 线}：我们先把 NTT、内积等做成向量探针；再把 Alg.13 拆成 $\hat{A}$、$s,e$、device 全链，逐段 PASS；最后集成晋级 stable。
  \item \textbf{KEM 线}：我们语义上调用 Alg.13（步骤 1 即 \texttt{K-PKE.KeyGen}）；设备上在 Launch-2 尾部追加 Alg.16 步骤 2--3（$H$、拼 $\mathit{dk}$）；编排上保持 2-launch；验收上 correctness 作 oracle，device 去 vendor 后晋级 stable。
\end{itemize}

\paragraph{设备侧 Launch 与 FIPS 行的对应（本仓实现）}
{\small
\begin{center}
\begin{tabular}{@{}L{2.0cm}L{5.5cm}L{2.2cm}L{4.0cm}@{}}
  \toprule
  Launch & 设备段 & FIPS 行 & 要点 \\
  \midrule
  L1 & \texttt{f203\_keygen\_prep} & Alg.13 步骤 1--15 & 双 AIV；$\hat{A}$、$\hat{s}$、$\hat{e}$ \\
  L2 & \texttt{mmad\_custom} + 尾缝 & Alg.13 步骤 16--21 & NTT、内积、ByteEncode \\
  L2 尾 & \texttt{KemKgTailFused} & Alg.16 步骤 2--3 & $z$、$H(\mathit{ek})$、写 $\mathit{dk}$ \\
  \bottomrule
\end{tabular}
\end{center}
}

若用一句话概括这次成功：
\textbf{我们把 FIPS 的三层算法，映射成「已证 PKE 链 + 一条单独想清楚的 KEM 尾缝」}，
而不是在 KEM 任务里重写 NTT。

\paragraph{图二：FIPS 三层的数据依赖（图感）}
图一讲「仓库里怎样长出来」；图二只看标准数据流——
\textbf{重计算在 Alg.13}，$z$ \textbf{不进} Alg.13，只进 Alg.16 的 $\mathit{dk}$ 拼接；
$\mathit{ek}_{\mathrm{KEM}}$ 对 $\mathit{ek}_{\mathrm{PKE}}$ 是直通。

\begin{center}
\begin{tikzpicture}[
  x=1.85cm, y=1.12cm,
  var/.style={draw, ellipse, minimum width=2.1cm, minimum height=0.76cm,
              font=\small, inner sep=1pt},
  op/.style={draw, rectangle, rounded corners=2pt, align=center, font=\small,
             minimum width=2.85cm, minimum height=0.9cm},
  arr/.style={-{Stealth[length=2.2mm]}, thick},
  lbl/.style={font=\footnotesize, text=black}
]
  % 节点
  \node[op, fill=orange!10] (a19) at (1.8, 5.5) {Alg.19\\采样 $d$、$z$};
  \node[var] (d) at (0.5, 4.2) {$d$};
  \node[var] (z) at (3.9, 4.2) {$z$};
  \node[op, fill=blue!10] (a13) at (0.5, 2.7) {Alg.13\\K-PKE.KeyGen};
  \node[var] (dkpke) at (-0.7, 1.2) {$\mathit{dk}_{\mathrm{PKE}}$};
  \node[var] (ekpke) at (1.8, 1.2) {$\mathit{ek}_{\mathrm{PKE}}$};
  \node[op, fill=gray!12, minimum width=2.0cm] (H) at (1.8, -0.25) {$H(\mathit{ek})$};
  \node[op, fill=orange!12, minimum width=5.0cm] (a16) at (1.5, -1.7)
    {Alg.16 拼 $\mathit{dk}$\\（$\mathit{ek}{\equiv}\mathit{ek}_{\mathrm{PKE}}$）\\
     $\mathit{dk}_{\mathrm{PKE}}\Vert\mathit{ek}\Vert H(\mathit{ek})\Vert z$};
  \node[var, fill=green!10] (ek) at (-0.2, -3.35) {$\mathit{ek}_{\mathrm{KEM}}$};
  \node[var, fill=green!10] (dk) at (2.9, -3.35) {$\mathit{dk}_{\mathrm{KEM}}$};
  \node[lbl, align=left, anchor=west] at (4.4, 2.0) {
    右泳道：$z$ 旁路\\
    \textbf{不}进入 Alg.13\\
    仅进 Alg.16
  };
  \node[lbl, align=center] at (1.4, -4.1) {底部绿色 $=$ KEM 交付输出};

  % 边：无交叉（左直通 / 中拼装 / 右旁路）
  \draw[arr] (a19) -- (d);
  \draw[arr] (a19) -- (z);
  \draw[arr] (d) -- (a13);
  \draw[arr] (a13) -- (dkpke);
  \draw[arr] (a13) -- (ekpke);
  \draw[arr] (ekpke) -- (H);
  \draw[arr] (dkpke) -- (a16);
  \draw[arr] (H) -- (a16);
  \draw[arr] (z.south) -- ++(0,-4.6) -| (a16.east);
  % 直通：从 $\mathit{ek}_{\mathrm{PKE}}$ 左侧外绕至交付叶（拼装式中 $\mathit{ek}$ 已标注 ${\equiv}$）
  \draw[arr] (ekpke.west) -- ++(-1.55,0) -- ++(0,-4.05)
    -| node[lbl, left, pos=0.2]{直通} (ek);
  \draw[arr] (a16) -- (dk);
\end{tikzpicture}
\end{center}

读图时抓住三点：
\begin{enumerate}
  \item \textbf{PKE 主体}：几乎全部重计算落在 Alg.13；$d$ 进、$\mathit{ek}_{\mathrm{PKE}}$/$\mathit{dk}_{\mathrm{PKE}}$ 出。
  \item \textbf{KEM 尾缝}：Alg.16 把已有片段拼成 $\mathit{dk}_{\mathrm{KEM}}$；
        $\mathit{ek}_{\mathrm{KEM}}$ 对 $\mathit{ek}_{\mathrm{PKE}}$ 直通（图底左侧绿叶），不重做 NTT。
  \item \textbf{汇点而非树}：NTT、CBD、MMAD 等还不画进本图——它们也不专属于 KeyGen，Encrypt / Encaps 会复用；
        因此本图只是更大有向无环图里的一个\textbf{汇点}，不是一棵只往下长的树。
\end{enumerate}

\subsection{从这个例子出发，我们应当问什么？}

例子讲完，真正重要的是\textbf{问题}。
如果我们的研究课题是「怎样让 Agent 预研少幻觉、少撞墙」，
那么面对 KEM.KeyGen 这个故事，我们至少应当认真问下面四组问题。

\paragraph{第一组：标准与工程之间的缝隙}
\begin{enumerate}
  \item FIPS 只告诉我们\textbf{算什么}；它并没有告诉我们，在昇腾上应当\textbf{先验哪几块、按什么顺序拼}。
        我们凭什么在动手前就说「可以直接站在 stable PKE 上」——凭的是目录名，还是某种可写的证书？
  \item 探针列表（INDEX）像一张「能力菜单」。
        \textbf{菜单本身}是否等于\textbf{合法拼装说明书}？若不等，缺的那一页应当长什么样？
\end{enumerate}

\paragraph{第二组：「单点都对」之后，组合还对吗}
\begin{enumerate}
  \setcounter{enumi}{2}
  \item 当 NTT 探针 PASS、PKE KeyGen PASS 之后，\textbf{谁}来保证「2-launch 接缝」「Alg.16 尾拼 $\mathit{dk}$」也一定 PASS？
        我们能否把这种「组合不免费」写进理论，而不是每次靠工程师的经验？
  \item 若我们把 correctness 路线（更多 launch、更保守同步）与 device/stable 路线（更少 launch、吃到优化）
        都称为「正确」，\textbf{我们究竟在验收什么}——仅仅是输出字节，还是也要验收「沿已证能力生长」？
\end{enumerate}

\paragraph{第三组：Agent 若只知道算法和菜单，会发生什么}
\begin{enumerate}
  \setcounter{enumi}{4}
  \item 若我们只给 Agent 一段 Alg.19 伪代码 + INDEX 里有哪些探针 PASS，
        \textbf{但不给}确定的拼装路线，搜索空间会有多大？
        我们亲眼见过：过程高度混乱、大量回退、极耗 token，最后只留下「能跑对」的 correctness 标本。
  \item 我们能否把搜索空间收束为「只允许沿已证书边生长」，同时又\textbf{不}把工程优化堵死？
        收束的规则应当写在图里，还是写在某种「合法句子」里？
\end{enumerate}

\paragraph{第四组：怎样把一次成功变成可复用的方法}
\begin{enumerate}
  \setcounter{enumi}{6}
  \item 这次 KEM.KeyGen 的成功，能否被复述成一张\textbf{闭包表}：
        依赖谁、缺了谁、哪条边有证书、哪条边被禁止？
        若下一个人（或 Agent）做 Encaps / Decaps，我们能否\textbf{先交这张表再写码}？
  \item 当上游改了 tiling 或同步壳，下游哪些节点应当自动标「脏」？
        我们是否需要一种比「记得上次 PASS」更可靠的失效传播？
\end{enumerate}

\subsection{本工程现在能回答这四组问题吗？}

四组问题都很好，但「能回答」要分层次：
\textbf{工程实践里我们已有碎片}，\textbf{形式化方法论里我们尚未闭环}。
下表诚实对照（2026-07 专题讨论时的判断）。

{\small
\begin{longtable}{@{}L{2.4cm}L{4.2cm}L{4.8cm}L{3.2cm}@{}}
  \toprule
  问题组 & 本工程\textbf{已能}说明的 & \textbf{仍缺}的 & 判断 \\
  \midrule
  第一组 &
  INDEX、customspec、baseline-registry 已在约束引用来源 &
  机器可读的「证书对象」；菜单$\neq$拼装说明书的显式规则 &
  \textbf{部分能答} \\
  \midrule
  第二组 &
  我们有 seam 探针/集成探针实践；correctness vs device 在 customspec 有对照 &
  「组合不免费」未写成统一理论；验收口径仍以 I/O 为主 &
  \textbf{部分能答} \\
  \midrule
  第三组 &
  correctness 历史是\textbf{反面教材}（高成本、不可控） &
  正向的「合法派生」门禁尚未工具化；Agent 仍可能绕开 &
  \textbf{能答负例；正例在建} \\
  \midrule
  第四组 &
  单任务可手写 INTEGRATION\_PLAN、闭包直觉 &
  统一闭包表 schema；dirty 传播无自动机制 &
  \textbf{部分能答} \\
  \bottomrule
\end{longtable}
}

\paragraph{逐组一句话}
\begin{itemize}
  \item \textbf{第一组}：我们能指着 stable PKE 说「KEM 应站在这里」，但还缺一张人人（含 Agent）都认的\textbf{证书页}。
  \item \textbf{第二组}：我们知道接缝要单独验（例如 2-launch、尾缝），但还没把这条经验收成\textbf{可检查的规则}。
  \item \textbf{第三组}：我们能用 correctness 证明「只给菜单会爆炸」；方法论要证明的是「给规则后能收敛」——这正是本专题要做的事。
  \item \textbf{第四组}：我们能复盘 KEM.KeyGen 的故事，但 Encaps/Decaps 还不能\textbf{先交闭包表再写码}作为硬门禁。
\end{itemize}

因此：\textbf{本工程能支撑这四组问题的提出与部分实证，但不能宣称方法论已经成立}。
后文的工作，就是把表中「仍缺的」一列逐项补上，并用 Encaps / Decaps 做检验。

\paragraph{这些问题的共同指向}

上面这些问题，表面上各不相同，但它们指向同一条研究主线：

\begin{quote}
我们能否把「已验证能力 + 合法组合 + 动态证书」整理成一套\textbf{可教、可检查、可对照实验}的形式方法——\\
让 Agent 的工作从「在菜单里瞎拼」变成「在自动机上走合法派生」？
\end{quote}

后文先补\textbf{方法论所需的数学背景}（第~\ref{sec:mathbg}~章：图、闭包、语言、自动机），
再把这些对象\textbf{装进工程语境}并写成门禁（第~\ref{sec:math}--\ref{sec:gates}~章），
最后用 Encaps 复盘、Decaps 前瞻检验这套语言是否管用。
\textbf{本章与后文均不展开} PQC / 数论变换等密码代数——那些细节见 \texttt{docs/notes/}；本专题专注方法论。

% ============================================================
\section{方法论所需的数学背景}
\label{sec:mathbg}

第~\ref{sec:hook}~章给了工程故事与问题感。
本章只补一件事：让读者具备读后续章节所需的\textbf{形式工具}。
我们从最基础的集合与关系讲起，一层层升到有向图、DAG、依赖闭包、形式语言与派生、配置与转移。
全程用第~\ref{sec:hook}~章的 KEM.KeyGen 例子回指；\textbf{不}讲格、MLWE、NTT 代数。

\subsection{本章要建什么、不建什么}

\begin{center}
\begin{tabular}{@{}L{3.2cm}L{5.8cm}L{5.0cm}@{}}
  \toprule
  & 本章建设 & 本章明确不建设 \\
  \midrule
  对象 & 预研过程的数学：依赖、合法展开、证书演化 & 密码方案的正确性证明 \\
  例子 & 只用 §2 已出现的 KeyGen 故事当插图 & 新开 Encaps/Decaps 代数推导 \\
  深度 & 直觉 + 可操作定义，够用即可 & 完备性、可判定性、复杂度（后置） \\
  \bottomrule
\end{tabular}
\end{center}

读完本章，你应当能用自己的话说清三句话：
\begin{enumerate}
  \item \textbf{DAG} 回答「引理库里谁依赖谁」；
  \item \textbf{派生 / 语法树} 回答「这一次任务怎样展开成句子」；
  \item \textbf{配置与转移} 回答「证书如何增减、何时接受或拒绝」。
\end{enumerate}

\subsection{集合、关系与函数（极短底座）}

设 $V$ 是我们将要讨论的「能力」或「中间结果」的名字集合。
\textbf{二元关系} $R\subseteq V\times V$ 就是「有序对的集合」：
若 $(u,v)\in R$，我们读作「从 $u$ 到 $v$ 有一条关系」。

\textbf{函数} $f:A\to B$ 是特殊关系：每个 $a\in A$ 恰对应一个 $f(a)\in B$。
后文里，「证书状态」「目标改写」都会用到这种「输入确定、输出唯一」的直觉；
暂时只需记住：\textbf{关系可以一对多，函数不允许一个输入两个结果}。

\paragraph{工程回指}
§2 图二里的箭头，就是 $V$ 上的一种关系：
「$d$ 参与 Alg.13」「$\mathit{ek}_{\mathrm{PKE}}$ 经直通得到 $\mathit{ek}_{\mathrm{KEM}}$」等。

\subsection{有向图}

有向图 $G=(V,E)$ 由\textbf{顶点集} $V$ 与\textbf{有向边集} $E\subseteq V\times V$ 组成。
一条边 $(u,v)\in E$ 画成 $u\to v$，读作「$u$ 指向 $v$」（方向有含义，不可随意反过来）。

\begin{itemize}
  \item \textbf{出边 / 入边}：从 $u$ 出发的边、进入 $u$ 的边；
  \item \textbf{路径}：顶点序列 $v_0,v_1,\ldots,v_k$，使得每步 $(v_i,v_{i+1})\in E$；
  \item \textbf{环}：起点与终点相同、且至少含一条边的路径。
\end{itemize}

若图中存在环，就可能出现「$A$ 依赖 $B$，同时 $B$ 又依赖 $A$」——
在「先验证下层再拼上层」的预研叙事里，这种环通常表示建模出错或任务切分不清。

\paragraph{工程回指}
§2 图一（探路）与图二（数据依赖）都是有向图的草图。
图二强调：$\mathit{dk}_{\mathrm{KEM}}$ 是多条入边的\textbf{汇合点}，不是一棵只向下分叉的树。

\subsection{DAG：有向无环图}

若有向图不含有向环，则称它为\textbf{有向无环图}（Directed Acyclic Graph，DAG）。

DAG 上一定存在\textbf{拓扑序}：把顶点排成一列 $v_1,\ldots,v_n$，
使得每条边都从较前的顶点指向较后的顶点。
直观说法：\textbf{可以按「先决条件在前」的顺序处理所有节点}。

\begin{center}
\begin{tikzpicture}[
  x=2.2cm, y=1.1cm,
  every node/.style={font=\small},
  box/.style={draw, rounded corners=2pt, align=center, minimum width=2.0cm, minimum height=0.7cm},
  arr/.style={-{Stealth[length=2mm]}, thick}
]
  \node[box, fill=blue!8] (p0) at (0,2) {P0 平台事实};
  \node[box, fill=blue!8] (p1) at (0,1) {P1 原语};
  \node[box, fill=blue!10] (p2) at (0,0) {P2 PKE.KeyGen};
  \node[box, fill=orange!12] (p3) at (0,-1) {P3 KEM.KeyGen};
  \draw[arr] (p0) -- (p1) -- (p2) -- (p3);
  \node[align=left, anchor=west] at (1.4,0.5) {示意：一条合法依赖链\\（真实图会有多父、多子）};
\end{tikzpicture}
\end{center}

\paragraph{为何预研爱用 DAG}
「先有证书的能力，再引用它拼更高层」天然要求依赖关系无环。
若出现环，要么把两个节点拆开，要么承认建模还没说清「谁先成立」。

\paragraph{重要区分（后文反复用）}
DAG 描述的是\textbf{引理库里的静态依赖投影}——「谁可以依赖谁」。
它\textbf{不}自动告诉你：某一次做 KEM.KeyGen 时，Agent 实际走了哪条展开路径。
那是下一节「派生」的事。

\subsection{依赖、前驱与闭包}

固定一张能力 DAG $G=(V,E)$。
对目标 $T\in V$，定义：

\begin{itemize}
  \item \textbf{直接依赖}：$\mathrm{Pred}(T)=\{u\mid (u,T)\in E\}$（画进 $T$ 的直接前驱）；
  \item \textbf{依赖闭包} $\mathrm{Dep}^*(T)$：从任意顶点出发、沿有向边能走到 $T$ 的所有顶点
        （含 $T$ 自己，或按约定不含——后文写闭包表时须声明）；
  \item \textbf{缺项}：闭包里尚未获得证书的顶点集合。
\end{itemize}

直觉算法（不必写成最优代码）：从 $T$ 出发\textbf{沿边反向}走，凡能走到的都收进闭包；
再与「当前已证书集合 $\Gamma$」做差，得到缺项列表。

\paragraph{工程回指}
做 KEM.KeyGen 之前，「闭包」至少应包含：stable PKE.KeyGen、Alg.16 尾缝相关能力、以及它们继续向下的 P1/P0 依赖。
§2 的资产清单，就是闭包表的工程雏形——只是当时还没叫这个名字。

\paragraph{与 §2 四组问题的扣合}
「菜单上有一排绿探针」$\neq$「闭包已算清且无缺项」。
INDEX 更像顶点菜单；\textbf{闭包}才回答「为了 $T$，我最少必须先有谁」。

\subsection{三个容易混的东西：DAG、派生、过程}

在进入形式语言之前，先把三个层面钉死：

\begin{center}
\begin{tabular}{@{}L{2.8cm}L{5.6cm}L{5.4cm}@{}}
  \toprule
  概念 & 回答的问题 & 不回答的问题 \\
  \midrule
  能力 DAG & 库里谁依赖谁？ & 某次任务一步步怎么展开？ \\
  派生 / 语法树 & 这次任务的展开句子是什么？ & 全局所有共享边的全貌？ \\
  配置与转移 & 证书如何增减？何时接受/拒绝？ & 输出字节是否符合 FIPS？ \\
  \bottomrule
\end{tabular}
\end{center}

\noindent
一句话：\textbf{DAG 是库；派生是一次作业；转移是作业过程中的状态演化。}
三者一起，才接近「预研怎样合法地做完」。

\subsection{字母表、字与语言}

\textbf{字母表} $\Sigma$ 是有限符号集。
$\Sigma$ 上的\textbf{字}（字符串）是有限序列 $w=a_1a_2\cdots a_n$（$a_i\in\Sigma$）；
空字记作 $\varepsilon$。
全体字的集合记作 $\Sigma^*$。

\textbf{语言} $L$ 就是 $\Sigma^*$ 的一个子集：$L\subseteq\Sigma^*$。
说「$w$ 合法」，就是说 $w\in L$。

\paragraph{方法论里怎么取 $\Sigma$（直觉）}
把「调用一个已证书能力」「做一个已声明的组合动作（如接一条接缝、加一次 launch 边界）」当成字母。
一条\textbf{实现计划}或\textbf{展开轨迹}就是一个字 $w\in\Sigma^*$；
全体合法计划构成语言 $L$。

\paragraph{工程回指}
「先做 stable PKE，再做 Alg.16 尾缝，再验收 KEM I/O」可以看成一个很短的字。
「在 KEM 任务里从零发明一套未证书的同步壳」则应落在 $L$ 之外——哪怕最后字节碰巧对了。

\subsection{产生式与派生（语法树）}

形式语言常用\textbf{产生式}描述「怎样从目标长出句子」。
直觉写法：
\[
T \;\Rightarrow\; u_1\, u_2\, \cdots\, u_k
\]
读作：目标 $T$ 可以被改写成子目标（或已证书调用）序列 $u_1\cdots u_k$，
\textbf{当且仅当}依赖与接缝已在图上声明、且改写规则允许。

从目标符号出发，反复使用产生式，得到一个字 $w$——这个过程叫\textbf{派生}。
派生的树形记录就是\textbf{语法树}：根是本次任务目标，叶子是已证书的原子步骤。

\begin{center}
\begin{tikzpicture}[
  every node/.style={font=\small, align=center},
  arr/.style={-{Stealth[length=2mm]}}
]
  \node (T) {KEM.KeyGen};
  \node[below left=8mm and 8mm of T] (P) {PKE.KeyGen\\（已证书）};
  \node[below right=8mm and 8mm of T] (S) {Alg.16 尾缝\\（接缝节点）};
  \draw[arr] (T) -- (P);
  \draw[arr] (T) -- (S);
  \node[below=5mm of P, font=\footnotesize] {继续展开到 P1/P0\ldots};
  \node[below=5mm of S, font=\footnotesize] {须单独获证};
\end{tikzpicture}
\end{center}

\paragraph{与 DAG 的再对照}
\begin{itemize}
  \item DAG 上可能有许多条边指向同一节点（多任务共享 NTT 等）；
  \item 某次派生树只选取\textbf{这一次任务需要的}展开，是 DAG 的一棵「使用痕迹」，不是整库重画。
\end{itemize}

\paragraph{成员判定（门禁雏形）}
Agent 提交一张「闭包表 + 展开计划」时，检查的是：
是否存在从目标到该计划的合法派生，且中途不出现禁字（见下）。
这就是后文「先交闭包表，再写码」的数学说法。

\subsection{禁字、否定与「组合不免费」}

语言不仅由「能写什么」定义，也由「不能写什么」定义。

\begin{itemize}
  \item \textbf{禁字 / 否定边}：显式禁止出现的片段——例如引用已冻结路线、使用未登记计算块、在默认生产路径上走 Host 密码胶水等。
        仓库里的 \texttt{frozen/} + \texttt{FROZEN.md}，就是否定边的工程文书。
  \item \textbf{组合不免费}：若 $A$、$B$ 各自对应合法片段，拼接 $A\!\circ\!B$ \textbf{不必}自动合法。
        往往需要一个新的\textbf{接缝节点}（或新产生式）并单独获证。
        §2 的「2-launch 接缝」「Alg.16 拼 $\mathit{dk}$」正是此类对象。
\end{itemize}

因此：$L$ 不是「把绿菜单任意排列组合」。
\textbf{菜单给出字母候选；产生式与否定边共同给出合法句子。}

\subsection{配置与转移（自动机直觉）}

自动机讨论的是：系统处于什么\textbf{配置}（状态），在什么规则下可以\textbf{转移}到下一配置，
以及哪些配置算\textbf{接受}或\textbf{拒绝}。

对本方法论，一种够用的配置写法是三元组
\[
q = (\Gamma,\, F,\, S),
\]
其中：
\begin{itemize}
  \item $\Gamma$：当前已证书能力集合（引理库的「已绿」部分）；
  \item $F$：前沿目标（还在展开/等待获证的任务）；
  \item $S$：接缝与约束（含否定边、I/O 契约等）。
\end{itemize}

典型转移（名字可后改，含义先固定）包括：
\begin{center}
\begin{tabular}{@{}L{2.6cm}L{11cm}@{}}
  \toprule
  转移 & 直觉 \\
  \midrule
  $\mathsf{Expand}$ & 把前沿目标按产生式展开，写入缺项 \\
  $\mathsf{Probe}$ & 对缺项做探针/预研，寻求证据 \\
  $\mathsf{Certify}$ & 验收通过后把节点写入 $\Gamma$ \\
  $\mathsf{Compose}$ & 声明并认证一条接缝（组合不免费的落点） \\
  $\mathsf{Dirty}$ & 上游变更导致下游证书失效（$\Gamma$ 可收缩） \\
  $\mathsf{Freeze}$ & 写入否定边，永久禁止某条路线 \\
  \bottomrule
\end{tabular}
\end{center}

\begin{itemize}
  \item \textbf{接受}：目标 $T$ 已进入 $\Gamma$，且生产 I/O 契约满足；
  \item \textbf{拒绝}：出现未证书边、踩中否定边、或跳层发明被门禁拦住。
\end{itemize}

\paragraph{工作定义（本章收束）}
\begin{quote}
预研开发的\textbf{合法性} $\approx$ 在已证书能力生成的语言 $L$ 上，
由配置转移搜索一条接受路径。\\
\textbf{幻觉} $\approx$ 非法转移，或画出/执行了无证书边。\\
FIPS、liboqs、golden 仍是语义 oracle——它们回答「算得对不对」；
形式方法回答「砖允许从哪来、允许怎样砌」。
\end{quote}

\subsection{非单调：证书为什么会「变脏」}

古典证明系统常被想象成「引理只增不减」。
工程引理库不是这样：上游 iface、tiling 或同步壳一变，下游「曾经 PASS」可能失效。
我们用状态 $\mathsf{dirty}$ 标记这种失效；对应转移是 $\mathsf{Dirty}$。

这意味着 $\Gamma$ \textbf{可收缩}。
后文门禁必须承认非单调，否则 Agent 会把过期证书当成永真公理。

\subsection{用 §2 四组问题做一次「翻译练习」}

\begin{center}
\begin{tabular}{@{}L{2.4cm}L{11.2cm}@{}}
  \toprule
  §2 问题组 & 用本章语言重述（一句话） \\
  \midrule
  第一组 & 标准给出语义 oracle；合法砌法由 DAG + 产生式 + 否定边给出，不是由目录名暗示。 \\
  第二组 & $A,B\in L$ 推不出 $A\!\circ\!B\in L$；接缝须 $\mathsf{Compose}$ 并获证。 \\
  第三组 & 只给伪代码 + 菜单 $\approx$ 只给 $\Sigma$ 的一部分，未给 $L$ 的成员判定；搜索空间爆炸。 \\
  第四组 & 一次成功应沉淀为闭包表 + 派生痕迹 +（必要时）新节点入 $\Gamma$；上游变更触发 $\mathsf{Dirty}$。 \\
  \bottomrule
\end{tabular}
\end{center}

\subsection{本章小结与下一章预告}

本章建立了方法论的最小数学大厦：
\textbf{有向图 / DAG $\to$ 依赖闭包 $\to$ 语言与派生 $\to$ 配置转移 $\to$ 证书非单调}。
刻意未做：密码代数、完备性证明、具体文件格式。

下一章（第~\ref{sec:math}~章）把这些对象\textbf{装进本仓库的工程字段}：
节点带哪些属性、边分哪几类、证书状态怎样命名——仍然是草案，可随复盘修改。

% ============================================================
\section{从数学到工程对象：节点、边与证书}
\label{sec:math}

上一章给了一般语言；本章把它\textbf{具象化}成预研图上的可填写字段。
若与第~\ref{sec:mathbg}~章用语冲突，以「工程可执行」优先，并回写修订记录。

\subsection{三个层面（复习）}
\begin{center}
\begin{tabular}{@{}L{3.2cm}L{5.5cm}L{5.3cm}@{}}
  \toprule
  概念 & 回答的问题 & 不回答的问题 \\
  \midrule
  能力 DAG & 引理库里谁依赖谁？ & 某次任务如何一步步展开？ \\
  语法树 / 派生 & 这次任务的展开句子是什么？ & 全局所有共享边的全貌？ \\
  配置自动机 & 证书如何增减？何时接受/拒绝？ & 代数是否符合 FIPS？ \\
  \bottomrule
\end{tabular}
\end{center}

\subsection{节点（已验证能力）}
节点 $v$ 建议携带：
\begin{itemize}
  \item \texttt{id}：稳定名（如 \texttt{pke.keygen.k4}）；
  \item \texttt{layer}：P0--P3（方法论分层，不等于目录树一一对应）；
  \item \texttt{iface}：I/O、dtype、形状/布局、同步约定——\textbf{可复用的是 iface，不是源码}；
  \item \texttt{cert}：路径 + CPU+SIM（+KAT）+ 版本锚点；
  \item \texttt{forbid}：否定约束；
  \item \texttt{repo}：权威目录锚点。
\end{itemize}

\paragraph{分层语言（P0--P3）}
\begin{center}
\begin{tabular}{@{}ll@{}}
  \toprule
  层 & 含义（本仓库语境） \\
  \midrule
  P3 & 顶层算子 / 算法入口（如 KEM.KeyGen、KEM.Encaps） \\
  P2 & 算法构件（如 PKE.KeyGen / Encrypt / Decrypt 全链） \\
  P1 & 领域原语（NTT、basemul、CBD、ByteEncode、SHAKE\ldots）——\textbf{作能力名}，本章不讲其代数 \\
  P0 & 平台公理（DataCopy、向量算子、MMAD+tiling、TPipe/TQue 同步壳\ldots） \\
  \bottomrule
\end{tabular}
\end{center}

\subsection{边的三种类型}
\begin{enumerate}
  \item \textbf{语义依赖}：$u$ 的正确性需要 $v$ 的 iface（如「NTT 能力」$\to$ PKE.KeyGen）；
  \item \textbf{接缝（compose）}：$A$、$B$ 各自有证，$A\!\circ\!B$ \textbf{仍须}单独成节点获证
        （「组合不免费」）；
  \item \textbf{否定边}：显式禁止（frozen 模板、未登记计算块、Host 默认密码路径等）。
\end{enumerate}

\subsection{证书状态（含非单调）}
\[
\mathsf{unproven}\to\mathsf{probe}\to\mathsf{lemma}\to\mathsf{deliverable},
\quad\text{以及}\quad
\mathsf{dirty}.
\]
\texttt{dirty} 表示上游 iface/tiling 变更后下游证书失效——引理库 $\Gamma$ \textbf{可收缩}。

\subsection{仓库机制与数学对象的对照（直觉）}
\begin{itemize}
  \item 活跃 \texttt{INDEX.md}：当前可引用的能力菜单（$\Sigma$ 的工程投影，\textbf{还不是} $L$）；
  \item \texttt{frozen/} + \texttt{FROZEN.md}：否定边；
  \item \texttt{baseline-registry}：golden 允许调用的已验证计算块；
  \item \texttt{customspec}：examples 写码的规格门禁（派生计划的文书形态之一）。
\end{itemize}

\subsection{立论（承接第~\ref{sec:mathbg}~章）}
\begin{quote}
预研开发的合法性 $\approx$ 在已证书能力生成的语言 $L$ 上，
由配置自动机搜索接受路径；\\
\textbf{幻觉} $\approx$ 非法转移 / 无证书边。\\
语义 oracle（FIPS/liboqs/golden）管「算对」；形式方法管「从哪砌砖」。
\end{quote}

% ============================================================
\section{门禁与工作流：Agent 被允许做什么}
\label{sec:gates}

\subsection{合法性规则（v0）}
接到目标 $T$ 时，只允许：
\begin{enumerate}
  \item 展开依赖闭包（语义边 + 接缝边）；
  \item 分类：已有 lemma/deliverable；缺项；仅否定边；
  \item 缺项 $\Rightarrow$ 先探针/exp 并获证，再引用；禁止在 $T$ 任务里顺便发明 P0/P1；
  \item 新接缝 $=$ 新节点，须认证；
  \item 验收看 I/O 与 golden/KAT，不看「与参考源码同构」；
  \item 自包含：挂 iface，不跨目录偷源码（仓库既有例外除外）；图边 $\neq$ \texttt{\#include}。
\end{enumerate}

\subsection{四步工作流}
\begin{lstlisting}
1. 目标声明  T、I/O、验收级别（probe / lemma / deliverable）
2. 闭包展开  P0--P3；ok / missing / dirty；接缝清单
3. 最小补洞  只补 missing 与新接缝；forbid 写入否定边
4. 闭环写回  新节点入图；失败则改规则或 Freeze，不默认可复用
\end{lstlisting}
\textbf{先交闭包表，再写码。} customspec / STATUS / registry 是文书形态，不是替代闭包表。

\subsection{域无关性}
抽出密码词之后，对象仍是：终端（axiom）、非终端（goal）、产生式（compose）、证书（judgment）。
同一套问题适用于非密码 AscendC 预研；ML-KEM 只是 oracle 清晰、接缝丰富的\textbf{压力测试床}。

% ============================================================
\section{复盘：理论如何对照 \texttt{kem.keygen} 与 \texttt{kem.encaps}}
\label{sec:replay}

\subsection{复盘的目的}
不是证明「我们早就在按理论做」（事后贴标签无价值），而是：
\begin{enumerate}
  \item 检查分层与接缝是否\textbf{可被理论语言说清}；
  \item 找出偏离（例如曾依赖 correctness vendor / frozen 模板）并记为\textbf{方法论债}或\textbf{否定边}；
  \item 总结「数学原理如何被具象化为目录、launch、证书」。
\end{enumerate}

\subsection{\texttt{kem.keygen} 校准问题单}
\begin{itemize}
  \item P3 是否只依赖已证书 P2，而非在 KeyGen 任务内重写 NTT？
  \item 2-launch、Alg.16 尾、\texttt{dk} 展开等接缝是否显式？
  \item baseline-registry 中的 golden 块是否都在图上有 lemma？
  \item 否定边是否写入（Host 默认禁密码、禁把 liboqs 当 AscendC 规格）？
\end{itemize}
仓库锚点（查阅用）：
\texttt{examples/stable/stable-fips203-mlkem-kem-keygen-k4/}，
\texttt{docs/specs/fips203-mlkem1024-kem-keygen-baseline-registry.md}，
及相关 \texttt{pass-fix-*} device 探针。

\subsection{\texttt{kem.encaps} 复盘要点（讨论共识）}
\begin{itemize}
  \item 合法路径应挂已证 \texttt{pke.encrypt} 与哈希/G 等 lemma，而不是以 correctness 内 vendor 为实现模板；
  \item 与 stable Encrypt 的编译期/接缝关系必须写进闭包，而不是「觉得复用了」；
  \item 若实现史存在混乱探索，复盘文应诚实记录——那是对照组与方法论的动机，不是羞耻。
\end{itemize}

\subsection{期望产出的「具象化对照表」（后续填空）}
\begin{center}
\begin{tabular}{@{}L{4cm}L{5cm}L{5cm}@{}}
  \toprule
  理论对象 & 工程具象 & 证书形态 \\
  \midrule
  节点 $v$ & 探针/exp/stable 目录 & STATUS + run.sh 对拍 \\
  语义边 & 「必须先有某某能力」 & INDEX / customspec 引用 \\
  接缝节点 & 多 launch / 融合边界 & 集成探针单独 PASS \\
  否定边 & frozen / forbid 条文 & FROZEN.md、Rule \\
  字 $w\in L$ & 一次任务的闭包表+计划 & 评审通过再写码 \\
  \bottomrule
\end{tabular}
\end{center}

% ============================================================
\section{前瞻：用理论指导尚未完成的 \texttt{kem.decaps}}
\label{sec:forward}

\subsection{为什么 Decaps 适合做试金石}
Decaps 同时依赖 Encrypt、Decrypt、KeyGen I/O 与 FO 失败路径，是典型的\textbf{多父 DAG}；
且若理论只能解释已完成故事、不能生成待做清单，则方法论停在修辞。

\subsection{前瞻时 Agent（或人类）应交付的最小闭包}
\begin{enumerate}
  \item 目标 I/O 与验收级别；
  \item 依赖节点清单（标 ok / missing / dirty）；
  \item 接缝清单（含对称失败 / 重加密比对等是否独立成节点）；
  \item 否定边（不可引用的 correctness vendor / frozen 实现）；
  \item 最小补洞顺序（先探针哪条缝）。
\end{enumerate}

\subsection{成功标准（阶段性）}
\begin{itemize}
  \item \textbf{弱成功}：闭包表可执行；偏差归因到缺边/假边/脏证书，并回写本章理论；
  \item \textbf{强成功}：在门禁下完成 device/stable 级实现，且过程指标优于 correctness 探索史
        （见下一章比较轴）。
\end{itemize}

% ============================================================
\section{对照组：\texttt{correctness} 标本的意义}
\label{sec:correctness}

\subsection{为何保留 correctness}
\texttt{*-correctness-*} 在\textbf{完全不考虑优化}的前提下，只追求可执行性与结果正确：
多 kernel launch、多搬运与同步——这是「一定能搭出正确积木」的保守拼法，
性能不可用，也几乎吃不进后续 P0/P1 工程优化。

它还记录了一段真实的 Agent 过程史：故意不给确定性拼装路线时，
探索成本极高（失败、卡死、回退；曾短期内耗尽大量 token），最终仅得到「结果正确」的代码。

\subsection{在方法论中的角色}
\begin{center}
\begin{tabular}{@{}L{3cm}L{11cm}@{}}
  \toprule
  角色 & 含义 \\
  \midrule
  正确性锚 & 下界存在性：积木可搭、I/O 可对 \\
  结构反例 & 展示「只求对」会长成怎样的 launch/同步膨胀 \\
  过程对照 & 无确定路线的搜索成本 \\
  实验基线 & 与方法论 Agent、人工指导路径三方比较 \\
  \bottomrule
\end{tabular}
\end{center}

\textbf{纪律}：correctness 是\textbf{基线标本}，不是活跃实现模板；可读、可对拍、可对比，\textbf{禁止}当抄码来源（与 frozen 出门不带码一致）。

\subsection{建议比较轴}
\begin{enumerate}
  \item \textbf{过程}：token / 轮次 / 回退；是否先交闭包表；非法边是否被拦截；
  \item \textbf{结构}：launch 数、同步点、lemma 复用、接缝是否显式；
  \item \textbf{结果}：CPU+SIM；SIM tick 相对 correctness 的倍数。
\end{enumerate}
三档：A 旧式无路线探索 $\sim$ correctness；B 新方法论 Agent；C 人工指导 device/stable。
B 不必一次打赢 C；明显优于 A，即方法论第一关。

% ============================================================
\section{未决议事项、日程与维护}
\label{sec:open}

\subsection{刻意未决议}
\begin{itemize}
  \item 节点是否 $1$:$1$ 对应目录名；图的存储是 Markdown 表还是 YAML/JSON；
  \item 主叙事选工作流网 / Petri 网 / 属性文法 / 判断式中的哪一种；
  \item 何时写入 Rule/Skill；可判定性与复杂度是否单开理论线。
\end{itemize}

\subsection{近期填空顺序（与「特殊$\to$一般」一致）}
\begin{enumerate}
  \item 按复盘反馈修订第~\ref{sec:mathbg}~章（定义精度、例子密度）；把 KeyGen 合法派生轨迹写成一页纸；
  \item 第~\ref{sec:math}~章字段表与仓库文书（闭包表 schema）对齐；
  \item Encaps 理论轨迹 vs 实现轨迹对照表；
  \item Decaps 前瞻闭包；
  \item correctness 对照实验写成可执行协议；
  \item 结论稳定后再考虑迁 \texttt{docs/notes/}。
\end{enumerate}

\subsection{文档与分支约定}
\begin{itemize}
  \item 本文件为专题\textbf{唯一} TeX/PDF 工作副本（无日期后缀），持续修订；
  \item 讨论纪要：\texttt{形式语言DAG预研讨论纪要.md}（同样单文件刷新）；
  \item 工作分支：\texttt{research/formal-lang-dag}；未经用户明确指令\textbf{不合入} \texttt{main}；
  \item 不记录题外的知识产权保护策略讨论。
\end{itemize}

\subsection{编译}
\begin{lstlisting}
bash scripts/xelatex-clean.sh \
  docs/research/从已验证能力到合法派生-面向Agent预研的形式方法教材草案.tex
\end{lstlisting}

\end{document}
